% =========================================================
% Bounds / Suprema / Infima — Review Exercises
% Aggregated Across Core Analysis Texts
% =========================================================

\section*{Bounds, Suprema, and Infima — Exercise Review Set}

These exercises train the foundational completeness machinery of the real
numbers. Every problem in this collection focuses on one of the following
techniques:

\begin{center}
\texttt{upper-bound argument \quad least-upper-bound property \quad epsilon-approximation}\\
\texttt{supremum algebra \quad sequence bounds \quad diameter via supremum}
\end{center}

Mastery of these problems prepares the student for monotone convergence,
Bolzano--Weierstrass, and compactness arguments.

\bigskip

% =========================================================
% Exercise Stub Template
% =========================================================

\newcommand{\exstub}[4]{
\subsection*{Exercise: #1}

\textbf{Source:} #2

\textbf{Technique:} #3

\textbf{Task.} #4

\begin{remark*}[Work Area]
Write your proof or computation here.
\end{remark*}

\bigskip
}

% =========================================================
% Abbott
% =========================================================

\section*{Exercises from Abbott}

\exstub
{Abbott 1.2.6}
{Abbott, \textit{Understanding Analysis}}
{least-upper-bound property}
{Prove that the supremum of a set is unique.}
% =========================================================
% Exercise Metadata
% tech:unpack-definition
% dep:def:upper-bound
% dep:def:supremum
% =========================================================

\ExerciseMeta
  {Abbott}
  {Understanding Analysis}
  {Exercise 1.2.6}
  {PROOF}
  {tech:unpack-definition}
  {dep:def:upper-bound, dep:def:supremum}

\begin{exerciseproofrecord}

\begin{remark*}[Technique]
\TechniqueEntry{tech:unpack-definition}
\end{remark*}

\begin{remark*}[Dependencies]
\DependencyEntry{dep:def:upper-bound},
\DependencyEntry{dep:def:supremum}
\end{remark*}

\begin{remark*}[Task]
Prove that the supremum of a set, if it exists, is unique.
\end{remark*}

\begin{remark*}[Strategy]
Assume two suprema exist. Since each is the least upper bound, each must be
less than or equal to the other. Then conclude equality.
\end{remark*}

\begin{proof}
Assume \(S\subseteq \mathbb{R}\), and suppose \(s_0\) and \(s_1\) are both
suprema of \(S\).

\tagDU Since \(s_0=\sup S\), by definition \(s_0\) is an upper bound of \(S\),
and \(s_0\) is less than or equal to every upper bound of \(S\).

\tagDU Likewise, since \(s_1=\sup S\), \(s_1\) is an upper bound of \(S\),
and \(s_1\) is less than or equal to every upper bound of \(S\).

\tagDU Because \(s_1\) is an upper bound of \(S\), and \(s_0\) is the least
upper bound, it follows that
\[
s_0 \le s_1.
\]

\tagDU Similarly, because \(s_0\) is an upper bound of \(S\), and \(s_1\) is
the least upper bound, it follows that
\[
s_1 \le s_0.
\]

Therefore,
\[
s_0 \le s_1 \le s_0.
\]

Hence \(s_0=s_1\).

Thus the supremum of a set, if it exists, is unique.
\end{proof}

\begin{remark*}[Proof shape]
Direct proof. Assume two candidates satisfy the defining property of supremum,
then use the least-upper-bound property in both directions to obtain mutual
inequalities and conclude equality.
\end{remark*}

\end{exerciseproofrecord}
\exstub
{Abbott 1.2.7}
{Abbott, \textit{Understanding Analysis}}
{epsilon-approximation}
{Show that if $s = \sup A$, then for every $\varepsilon > 0$
there exists $a \in A$ with $s - \varepsilon < a \le s$.}
\begin{proof}
Assume \(s=\sup A\).

By definition of supremum, \(s\) is an upper bound for \(A\). Thus
\[
a \le s \qquad \text{for all } a\in A.
\]

Now let \(\varepsilon>0\) be arbitrary.

Because \(s-\varepsilon < s\) and \(s\) is the \emph{least} upper bound of \(A\), the number
\[
s-\varepsilon
\]
cannot be an upper bound of \(A\).

Therefore, by the negation of the definition of upper bound, there exists some \(a\in A\) such that
\[
a > s-\varepsilon.
\]

Since \(s\) is an upper bound for \(A\), we also have
\[
a \le s.
\]

Hence
\[
s-\varepsilon < a \le s.
\]

As \(\varepsilon>0\) was arbitrary, the result follows.
\end{proof}
\exstub
{Abbott 1.2.8}
{Abbott, \textit{Understanding Analysis}}
{least-upper-bound property}
{Let $A$ be bounded above. Show that $s=\sup A$ iff
(i) $s$ is an upper bound and
(ii) no smaller number is an upper bound.}

\exstub
{Abbott 2.3.1}
{Abbott, \textit{Understanding Analysis}}
{sequence bounds}
{Show that every convergent sequence is bounded.}

\exstub
{Abbott 2.3.3}
{Abbott, \textit{Understanding Analysis}}
{sequence bounds}
{Give an example of a bounded sequence that does not converge.}

% =========================================================
% Bartle
% =========================================================

\section*{Exercises from Bartle \& Sherbert}

\exstub
{Bartle 1.2.5}
{Bartle \& Sherbert, \textit{Introduction to Real Analysis}}
{least-upper-bound property}
{Prove that if $A$ is nonempty and bounded above,
then $\sup A$ exists.}

\exstub
{Bartle 1.2.6}
{Bartle \& Sherbert}
{supremum algebra}
{Show that $\sup(A+B) \le \sup A + \sup B$.}

\exstub
{Bartle 1.2.8}
{Bartle \& Sherbert}
{supremum algebra}
{Let $c>0$. Prove $\sup(cA)=c\sup(A)$.}

\exstub
{Bartle 1.2.9}
{Bartle \& Sherbert}
{supremum algebra}
{If $c<0$, prove $\sup(cA)=c\inf(A)$.}

\exstub
{Bartle 3.2.3}
{Bartle \& Sherbert}
{sequence bounds}
{Show that every bounded monotone sequence converges.}

% =========================================================
% Bruckner
% =========================================================

\section*{Exercises from Bruckner}

\exstub
{Bruckner 1.4.1}
{Bruckner, \textit{Real Analysis}}
{upper-bound argument}
{Show that every finite set of real numbers has a supremum
and an infimum.}

\exstub
{Bruckner 1.4.3}
{Bruckner}
{epsilon-approximation}
{Let $s=\sup A$. Construct a sequence $(a_n)$ in $A$
such that $a_n \to s$.}

\exstub
{Bruckner 1.4.5}
{Bruckner}
{sequence bounds}
{Show that if $(a_n)$ is bounded then
$\sup\{a_n : n \in \mathbb N\}$ exists.}

\exstub
{Bruckner 2.1.4}
{Bruckner}
{sequence bounds}
{If $a_n \to a$, prove the sequence is bounded.}

\exstub
{Bruckner 2.1.6}
{Bruckner}
{supremum algebra}
{Let $A,B$ be bounded sets.
Prove $\sup(A+B) \le \sup A + \sup B$.}

% =========================================================
% Tao
% =========================================================

\section*{Exercises from Tao}

\exstub
{Tao 6.3.1}
{Tao, \textit{Analysis I}}
{sequence bounds}
{Let $a_n = 1/n$. Determine $\sup\{a_n\}$ and $\inf\{a_n\}$.}

\exstub
{Tao 6.3.2}
{Tao}
{sequence bounds}
{Let $a_n=(-1)^n$. Find $\sup\{a_n\}$ and $\inf\{a_n\}$.}

\exstub
{Tao 6.3.3}
{Tao}
{sequence bounds}
{Let $a_n=n$. Determine the supremum and infimum of the sequence.}

% =========================================================
% Lebl
% =========================================================

\section*{Exercises from Lebl}

\exstub
{Lebl 1.3.6}
{Lebl, \textit{Basic Analysis I}}
{supremum algebra}
{Show that $\sup(x+A)=x+\sup(A)$.}

\exstub
{Lebl 1.3.7}
{Lebl}
{supremum algebra}
{Show that $\inf(x+A)=x+\inf(A)$.}

\exstub
{Lebl 1.3.8}
{Lebl}
{supremum algebra}
{Show that $\sup(xA)=x\sup(A)$ for $x>0$.}

% =========================================================
% Metric-space bounds
% =========================================================

\section*{Metric Space Bounds}

\exstub
{Topology 4.6.13}
{Point-Set Topology}
{diameter via supremum}
{Prove that $\operatorname{diam}(X)=\sup\{d(x,y):x,y\in X\}$.}

\exstub
{Topology 4.6.15}
{Point-Set Topology}
{sequence bounds}
{If $x_n \to x$ in a metric space,
show $\{x_n\}\cup\{x\}$ is bounded.}

\exstub
{Topology 4.6.12}
{Point-Set Topology}
{distance computation}
{Compute the distance from a point to a subset in $\mathbb R^2$.}

